% Caso de Uso: Consultar Reservaciones
% Sistema: Hotel Perros

%-------------------------------------- COMIENZA descripción del caso de uso.

\begin{UseCase}{CU16}{Consultar Reservaciones}{
    Permite al usuario visualizar el listado de las reservaciones registradas en el sistema y acceder al detalle específico de cada una.
}
    \UCitem{Versión}{\color{Gray}1.0}
    \UCitem{Autor}{\color{Gray}Axel}
    \UCitem{Supervisa}{\color{Gray}[Nombre del Supervisor]}
    \UCitem{Actor}{\hyperlink{Propietario}{Propietario} / \hyperlink{Administrador}{Administrador}}
    \UCitem{Propósito}{Conocer el estado, fechas y costos de las estancias programadas o pasadas.}
    \UCitem{Entradas}{Filtros de búsqueda (Opcional: por fecha, estado o mascota).}
    \UCitem{Origen}{Mouse / Pantalla táctil}
    \UCitem{Salidas}{Lista de reservaciones con información resumida (Mascota, Habitación, Fechas, Estado).}
    \UCitem{Destino}{Pantalla}
    \UCitem{Precondiciones}{
        \begin{itemize}
            \item El usuario debe haber iniciado sesión.
        \end{itemize}
    }
    \UCitem{Postcondiciones}{El sistema presenta la información solicitada filtrada según el rol del usuario.}
    \UCitem{Errores}{
        \begin{itemize}
            \item \textbf{E1}: No se encontraron reservaciones registradas.
            \item \textbf{E2}: Error de conexión al recuperar los datos.
        \end{itemize}
    }
    \UCitem{Tipo}{Primario}
\end{UseCase}

%--------------------------------------
% Trayectoria Principal
%--------------------------------------
\begin{UCtrayectoria}
    \UCpaso[\UCactor] Accede al módulo de "Reservaciones" desde el menú principal.
    \UCpaso El sistema verifica el rol del usuario autenticado.
    \UCpaso El sistema recupera las reservaciones correspondientes:
        \begin{itemize}
            \item Si es \textbf{Administrador}: Recupera todas las reservaciones del sistema.
            \item Si es \textbf{Propietario}: Recupera únicamente las reservaciones asociadas a su ID de usuario.
        \end{itemize}
    \UCpaso Muestra el listado de reservaciones ordenadas (ej. por fecha más próxima), indicando: Nombre de Mascota, Habitación, Fechas de Ingreso/Salida y Estado (Pendiente, Confirmada, Cancelada).
    \UCpaso[\UCactor] Selecciona una reservación específica haciendo clic en ella o en el botón \IUbutton{Ver Detalle}.
    \UCpaso El sistema solicita la información completa de la reservación seleccionada (ID).
    \UCpaso El sistema valida que la reservación pertenezca al usuario o que este sea administrador.
    \UCpaso Muestra la pantalla de "Detalle de Reservación" con la información desglosada:
        \begin{itemize}
            \item Datos generales (Mascota, Habitación, Fechas).
            \item Desglose financiero (Costo hospedaje, Costo servicios, Total).
            \item Lista de servicios adicionales contratados.
            \item Notas especiales.
        \end{itemize}
\end{UCtrayectoria}

%--------------------------------------
% Trayectorias Alternativas
%--------------------------------------

% Trayectoria A: Sin reservaciones
\begin{UCtrayectoriaA}{A}{No existen reservaciones registradas}
    \UCpaso El sistema busca en la base de datos y no encuentra registros asociados al usuario.
    \UCpaso Muestra el mensaje ``No tienes reservaciones registradas''.
    \UCpaso Muestra el botón \IUbutton{Nueva Reservación} invitando a crear una.
\end{UCtrayectoriaA}

% Trayectoria B: Acceso Denegado a Detalle
\begin{UCtrayectoriaA}{B}{Intento de acceso a reservación ajena (URL directa)}
    \UCpaso El usuario intenta acceder a una reservación específica mediante URL o historial.
    \UCpaso El sistema detecta que el ID de la reservación no pertenece al propietario actual (y no es admin).
    \UCpaso Retorna un error de permisos (403 Forbidden).
    \UCpaso Muestra el mensaje de error {\bf MSG-Error} ``No autorizado''.
    \UCpaso Redirige al listado general de reservaciones.
\end{UCtrayectoriaA}

%--------------------------------------
% Puntos de extensión
%--------------------------------------
\subsection{Puntos de extensión}
\UCExtenssionPoint{
    % Cuando:
    El usuario desea cancelar una reservación activa.
}{
    % Durante la región:
    Paso 8 (Visualización del Detalle).
}{
    % Casos de uso a los que extiende:
    \hyperlink{CU17}{CU17 Eliminar Reservación}.
}